\documentclass[a4paper,10pt]{article}
\usepackage[utf8]{inputenc}
\usepackage{geometry}
\usepackage{enumitem}
\usepackage[colorlinks=true, linkcolor=blue, urlcolor=blue, citecolor=green]{hyperref}
\usepackage{sectsty}

\geometry{margin=0.65in}
\setlist[itemize]{left=0pt, itemsep=0pt}
\subsubsectionfont{\normalfont}

\begin{document}
\pagestyle{empty}

\begin{center}
  \href{https://github.com/zguztorres/Bioinformatics-Portfolio}{Portafolio GitHub} \\[0.5em]
  {\LARGE \bfseries Zaira Guzmán} \\[0.5em]
  \href{mailto:zguztorres99@gmail.com}{zguztorres99@gmail.com} \quad | \quad \href{https://www.linkedin.com/in/zaira-guzman/}{LinkedIn} \quad | \quad Salamanca, España
\end{center}

\hrule
\vspace{4pt}

% Resumen Profesional
\subsection*{Resumen}
Biotecnóloga especializada en Bioinformática y Genómica Computacional con una sólida trayectoria en el análisis de datos biológicos y la biología sintética. Experta en Python, R, entornos Unix y procesamiento de datos ómicos, con experiencia demostrada en la limpieza de bases de datos complejas, creación de flujos de trabajo eficientes y evaluación de modelos de inteligencia artificial. Apasionada por el desarrollo de herramientas de código abierto y la aplicación de metodologías computacionales para impulsar la investigación científica y la sostenibilidad.

% Formación Profesional
\subsection*{Formación Profesional}

\textbf{Especialización en Bioinformática y Genómica Computacional} \hfill \textit{Octubre 2025 – Actualidad}
\begin{itemize}
    \item Universidad de Salamanca, España
\end{itemize}

\textbf{Ingeniería en Biotecnología} \hfill \textit{Finalizado en 2023}
\begin{itemize}
    \item Universidad Dr. José Matías Delgado, El Salvador
\end{itemize}


% Experiencia Laboral
\subsection*{Experiencia Laboral}

\subsubsection*{%
  \textbf{Science Specialist} en \href{https://www.linkedin.com/company/meridial-marketplace}{Meridial Marketplace} para \href{https://www.linkedin.com/company/invisible-technologies-inc-/}{Invisible Technologies},
  \hfill
  \textit{Diciembre 2024 -- Actualidad}
}
\begin{itemize}
  \item Evaluación comparativa (SxS) de modelos de IA generativa mediante diseño de prompts y análisis de rendimiento.
  \item Gestión de métricas operativas (AHT) y anotación de datos de precisión para el entrenamiento de modelos de lenguaje.
  \item Auditoría de interacciones de voz y recuperación de datos en tiempo real para evaluar la robustez de modelos LLM.
\end{itemize}

\subsubsection*{%
  \textbf{Becaria en Genómica y Bioinformática (Prácticas)} en \href{https://www.linkedin.com/company/genomics-research-center-genomac-hub/posts/?feedView=all}{Genomac Institute}, Virtual
  \hfill
  \textit{Agosto 2025 -- Enero 2026}
}
\begin{itemize}
  \item Diseño y optimización de péptidos terapéuticos mediante Machine Learning, utilizando modelos predictivos (CAMPR3 y Anti-CP) para identificar candidatos con alta actividad antimicrobiana y baja toxicidad.
  \item Descubrimiento de genes biosintéticos y metabolitos secundarios a través de análisis genómico, metagenómico y herramientas de predicción computacional (antiSMASH, Bagel4).
  \item Ejecución de protocolos de cribado virtual y acoplamiento molecular (Molecular Docking) para la validación de interacciones proteína-ligando.
\end{itemize}

\subsubsection*{%
  \textbf{Asistente Administrativo} en Arrocera La Palma, El Salvador
  \hfill
  \textit{Octubre 2023 -- Abril 2025}
}
\begin{itemize}
  \item Gestión administrativa, supervisión de estados financieros, control interno y coordinación de auditorías.
  \item Gestión de relaciones con clientes y control documental, garantizando el seguimiento preciso de cada cuenta y registro.
  \item Administración de datos de producción y gestión de residuos. 
\end{itemize}

\subsubsection*{%
  \textbf{Investigadora en el área de Agua, Saneamiento e Higiene} en \href{https://www.linkedin.com/company/stockholm-international-water-institute/posts/?feedView=all}{Instituto Internacional del Agua de Estocolmo SIWI (voluntariado)}, El Salvador
  \hfill
  \textit{Noviembre 2023 -- Junio 2024}
}
\begin{itemize}
  \item Ejecución de diagnósticos territoriales basados en datos (WASH) para identificar brechas críticas de acceso, saneamiento e higiene en comunidades vulnerables.
  \item Liderazgo de estrategias de incidencia y comunicación técnica ante autoridades gubernamentales y ONGs, transformando hallazgos científicos en hojas de ruta para la toma de decisiones.
  \item Diseño de soluciones sostenibles y programas de resiliencia comunitaria fundamentados en evidencia, orientados al empoderamiento local y la autogestión de recursos hídricos.
\end{itemize}

% Habilidades
\subsection*{Habilidades}
\begin{itemize}
  \item \textbf{Nuevas Tecnologías y Bioinformática:} Programación orientada al análisis de datos (Python, R, entornos Unix). Implementación de pipelines bioinformáticos en la nube aplicados a proteómica. Machine Learning en entornos multi-ómica. Manejo de bases de datos científicas.
  \item \textbf{Habilidades Transversales:} Gestión de proyectos (Metodologías ágiles), divulgación científica, redacción de informes técnicos, planificación, trabajo en equipos multidisciplinarios e inglés profesional.
\end{itemize}

% Formación Complementaria
\subsection*{Formación Complementaria y Congresos}

\textbf{Fundamentals of Data-Driven Precision Medicine for Diabetes} \hfill \textit{Diciembre 2025}
\begin{itemize}
    \item Stanford Department of Medicine
\end{itemize}

\textbf{Introducción a la Proteómica y sus aplicaciones BioMédicas} \hfill \textit{Diciembre 2025}
\begin{itemize}
    \item IBSAL, Salamanca
\end{itemize}

\textbf{Bioinformatics Course 201} \hfill \textit{Octubre 2025}
\begin{itemize}
    \item Genomac Institute Inc
\end{itemize}

\textbf{II Congreso Salvadoreño de Biotecnología y Biología Molecular (Ponente)} \hfill \textit{Junio 2024}
\begin{itemize}
    \item El Salvador
\end{itemize}

\textbf{Introduction to Data Analytics} \hfill \textit{Junio 2023}
\begin{itemize}
    \item IBM - Coursera
\end{itemize}

% Idiomas
\subsection*{Idiomas}
\begin{itemize}
  \item Español: Nativo
  \item Inglés: B2
\end{itemize}

\end{document}